\section{Soft decomposition}\label{sec:soft}

\begin{theorem}[S: residual decomposition]\label{thm:soft}
\formal{Decomposition.lean}{decomp_ineq_residual,decomp_ineq_soft}
Let $h:2^V\to\R$ be a polymatroid and let
$\Decomp(N,U,(H_1,\ldots,H_t))$. With $U_j$ as above,
\begin{align}
2\bigl(h(V)-h(U)\bigr)
&\le \sum_{v\notin U}h(v\mid U)
   +\sum_{j=1}^t\sum_{v\in H_j}h(v\mid N(v)\cup U_j)
\label{eq:soft-pieces}\\
&\le \sum_{v\notin U}h(v\mid U)
   +\sum_{v\notin U}h(v\mid N(v)\cup U).
\label{eq:soft-simple}
\end{align}
No recovery equality is assumed.
\end{theorem}
\begin{outline}
In a piece, work with $g=h_U$. A scan step with supplying context
$K=S\cup P_{b(v)}$ and outside parent $w$ has the exact identity
\[
g(v\mid K)
=I_g(w;K\cup\{v\})-I_g(w;K)+g(v\mid K\cup\{w\}).
\]
The last term is at most $h(v\mid N(v)\cup U)$ by the parent
containment condition. A step without an outside parent has the same
residual bound directly. If $v$ has already been scanned, its increment
is zero; repetitions are therefore not charged again. The nested
outside-parent potentials telescope, giving the piece bound
\[
2h(H\mid U)\le
\sum_{v\in H}h(v\mid U)
+\sum_{w\notin U\cup H} I_h(w;H\mid U)
+\sum_{v\in H}h(v\mid N(v)\cup U).
\]
The original elimination inequality supplies its first two terms;
\lean{totalDefect_le_residual} and \lean{piece_ineq_residual} keep the
residuals. Induct over pieces. The identity
$h(w\mid U\cup H)=h(w\mid U)-I_h(w;H\mid U)$ cancels the outside
potentials with the next piece's singleton terms. This proves
\eqref{eq:soft-pieces}. Conditional monotonicity and the disjoint
partition of $V\setminus U$ give \eqref{eq:soft-simple}.
\end{outline}

The retained contexts $U_j$ can reduce the residual bill. The simpler
form is useful when only a uniform estimate before the piece scan is
available. Both allow uncertified dependencies to be paid as remaining
uncertainty while keeping the same certified relation $N$.

\begin{corollary}[Common observations]\label{cor:conditional-soft}
\formal{Observations.lean}{conditional_soft,entropy_conditional_soft}
Let $X_v=W_{f(v)}$, let $Z$ be any common tuple of ambient variables,
and suppose $\Decomp(N,\varnothing,\mathcal H)$. Then
\begin{equation}
2H(X_V\mid Z)\le
\sum_v H(X_v\mid Z)+\sum_v H(X_v\mid X_{N(v)},Z).
\label{eq:conditional-soft}
\end{equation}
The law is arbitrary on a finite sample space and finite alphabet.
The same inequality holds for conditional ranks of any ambient
polymatroid under any coordinate map $f$.
\end{corollary}
\begin{outline}
Apply \cref{thm:soft} to the polymatroid
$S\mapsto h(f(S)\cup Z)-h(Z)$ with $U=\varnothing$.
The finite entropy theorem substitutes the existing
\leanfile{MotionDepthCodes.codeEntropy} instance.
\end{outline}
Thus common information may include nonlinear observations and
correlated auxiliary variables; independent input coordinates are not
required.

\begin{corollary}[Exact recovery]\label{cor:exact}
\formal{Decomposition.lean}{exact_decomp,exact_rateBound,exact_card_sq}
Under the hypotheses of \cref{thm:soft}, if $h(v\mid N(v))=0$ for
all $v$, the residuals vanish and
\[
2\bigl(h(V)-h(U)\bigr)\le\sum_{v\notin U}h(v\mid U).
\]
In particular a decomposition from $\varnothing$ gives the original
half-rate bound. For every storage code $C\subseteq\alpha^V$ over a
finite nonempty alphabet, it gives $|C|^2\le |\alpha|^{|V|}$.
Here a storage code means that agreement on $N(v)$ determines agreement
at $v$ for every pair of words in $C$.
\end{corollary}
\begin{outline}
Nonnegativity and conditional monotonicity force every residual to
zero. The original rate-to-cardinality theorem then supplies the
cardinality claim, including the empty-code case.
\end{outline}
This derives the earlier exact theorem directly from the extension.

\begin{corollary}[Relative Class S]\label{cor:class-s}
\formal{ClassS.lean}{classSOn_soft,classS_soft}
If $N$ satisfies relative Class S on $R\subseteq V$ and $h$ is any
polymatroid, then
\[
2\bigl(h(V)-h(V\setminus R)\bigr)
\le \sum_{v\in R}h(v\mid V\setminus R)
 +\sum_{v\in R}h(v\mid N(v)\cup(V\setminus R)).
\]
Taking $R=V$ gives the global soft inequality on Class S.
\end{corollary}
\begin{outline}
Use the tenth paper's relative decomposition certificate and apply
\cref{thm:soft} with $U=V\setminus R$.
\end{outline}
Existing Class S certificates therefore remain usable with explicit
nonzero recovery uncertainty; no additional graph classification is
claimed.
